\documentclass[11pt, a4paper]{report}

\usepackage[utf8]{inputenc}
\usepackage[T1]{fontenc}
\usepackage[french]{babel}

% Appel des packages dans le dossier preambule
\def\packagepath{../../../../../preambule}   % path du package princial

\usepackage{\packagepath/page_layout} % <--- Nouveau
\usepackage{\packagepath/config_info}
\usepackage{\packagepath/cadres_cours}
\usepackage{\packagepath/zones_reponse}
\usepackage{\packagepath/config_ds}
\usepackage{\packagepath/config_maths}

% --- PILOTAGE ---
%\def\visibleornot{visible}    % visible or invisible
\ifdefined\visibleornot\else\def\visibleornot{visible}\fi


\def\documentpath{./Sources_Latex}
\graphicspath{{./Images/}}


\def\ptsexoA{0}



\begin{document}

\def\level{Premiere}  % Classe à droite
\def\course{NSI}
\def\chaptername{Les dictionnaires} % Titre au centre
\def\docname{AP07~--~TD 1}        % Identifiant à gauche

\pagestyle{DS_LP}

\NomPrenom{}

\begin{exercice_cours}[Indiquer ce qui est affiché par les instructions suivantes:
]{}
    
\begin{CodePythontexAlt}[Largeur=1\linewidth,Centre,Lignes=false,PremLigne=1]{}
dico = {'nom': 'Dupont', 'age': 30, 'ville': 'Paris', 'enfants': ['Alice', 'Bob']}
\end{CodePythontexAlt}



\begin{minipage}{0.49\linewidth}
    \begin{tabular}{p{0.4\linewidth}p{0.4\linewidth}}
    \hline
    Instruction & Affichage \\    \hline
    \texttt{dico['nom']} & \reponseLigne[3.cm]{\texttt{Dupont}}  \\
    \texttt{dico['age']} & \reponseLigne[3.cm]{\texttt{30}}  \\
    \texttt{dico['ville']} & \reponseLigne[3.cm]{\texttt{Paris}}  \\
    \hline  
    \end{tabular}
\end{minipage}\hfill{}
\begin{minipage}{0.49\linewidth}
\begin{tabular}{p{0.4\linewidth}p{0.4\linewidth}}
    \hline
    Instruction & Affichage \\    \hline
    \texttt{len(dico)} & \reponseLigne[3.cm]{\texttt{4}}  \\
    \texttt{dico['enfants'][0]} & \reponseLigne[3.cm]{\texttt{Alice}}  \\
    \texttt{dico['enfants']} & \reponseLigne[3.cm]{\texttt{['Alice', 'Bob']}}  \\
   \hline  
    \end{tabular}

\end{minipage} 

\end{exercice_cours}

\begin{exercice_cours}[On considère le dictionnaire suivant: ]{}
    
    \begin{CodePythontexAlt}[Largeur=1\linewidth,Centre,Lignes=false,PremLigne=1]{}
dico = {'a': 1, 'b': 2, 'c': 3}
\end{CodePythontexAlt}

Ecrire les instructions permettant
    \begin{enumerate}
        \item d'ajouter une nouvelle entrée 'd' avec la valeur 4. \hfill\reponseLigne[4cm]{\texttt{dico['d'] = 4}}
        \item de modifier la valeur associée à la clé 'b' pour qu'elle soit égale à 5. \hfill\reponseLigne[4cm]{\texttt{dico['b'] = 5}}
        \item de récupérer la valeur associée à la clé 'c'. \hfill\reponseLigne[4cm]{\texttt{dico['c']}}
        \item de vérifier si la clé 'a' est présente dans le dictionnaire. \hfill\reponseLigne[4cm]{\texttt{'a' in dico.keys()}}
        \item de supprimer l'entrée associée à la clé 'a'. \hfill\reponseLigne[4cm]{\texttt{del dico['a']}}
    \end{enumerate}

\end{exercice_cours}

\begin{exercice_cours}[On considère la liste suivante: ]{}
    
\begin{CodePythontexAlt}[Largeur=1\linewidth,Centre,Lignes=false,PremLigne=1]{}
L = ['Alice', 'Bob', 'Charlie', 'David']
\end{CodePythontexAlt}

Ecrire une fonction \verb|longueur_mot(L)| qui renvoie un dictionnaire dont:
\begin{itemize}
    \item les clés sont les éléments de la liste L
    \item les valeurs sont les longueurs des éléments de la liste L
\end{itemize}


        \begin{CodePythontexAlt}[TaillePolice=\large, Largeur=1\linewidth,Centre,Lignes=false]{}
    def longueur_mot(L):
        
    





    assert longueur_mot(L) == {'Alice': 5, 'Bob': 3, 'Charlie': 7, 'David': 5}
        \end{CodePythontexAlt}


       
\end{exercice_cours}

\pagebreak

\begin{exercice_cours}[\quad]{}
    
Dans cet exercice, on gère une petite guilde de guerriers. Le dictionnaire guilde contient le nom des joueurs en clés et leur niveau en valeurs.

\begin{CodePythontexAlt}[Largeur=1\linewidth,Centre,Lignes=false,PremLigne=1]{}
guilde = {
    "Aethelgard": 42,
    "Slythe": 15,
    "Valeria": 38,
    "Grognar": 10,
    "Lyra": 25
}
\end{CodePythontexAlt}

\begin{enumerate}
    \item Ecrire une fonction \verb|liste_de_noms(guilde)| qui renvoie la liste des noms des joueurs. \hfill
    
    \begin{CodePythontexAlt}[TaillePolice=\large, Largeur=1\linewidth,Centre,Lignes=false]{}
def liste_de_noms(guilde):







assert liste_de_noms(guilde) == ["Aethelgard", "Slythe", "Valeria", 
                                 "Grognar", "Lyra"]
        
    \end{CodePythontexAlt}

    \item Ecrire une fonction \verb|niveau_moyen(guilde)| qui renvoie le niveau moyen des joueurs de la guilde. \hfill

    \begin{CodePythontexAlt}[TaillePolice=\large, Largeur=1\linewidth,Centre,Lignes=false]{}
def niveau_moyen(guilde):







assert niveau_moyen(guilde) == 26.0
    \end{CodePythontexAlt}

\item Ecrire une fonction \verb|listes_elite(guilde)| qui renvoie un dictionnaire contenant uniquement les joueurs de niveau supérieur ou égal à 30. \hfill

    \begin{CodePythontexAlt}[TaillePolice=\large, Largeur=1\linewidth,Centre,Lignes=false]{}
def listes_elite(guilde):











assert listes_elite(guilde) == {"Aethelgard": 42, "Valeria": 38}
    \end{CodePythontexAlt}
\end{enumerate}


       
\end{exercice_cours}


    

\end{document}